\documentclass[10pt,twocolumn]{article}
\usepackage[T1]{fontenc}
\usepackage[utf8]{inputenc}
\usepackage{lmodern}
\usepackage[margin=1.8cm,top=2.4cm,bottom=2cm,columnsep=0.6cm]{geometry}
\usepackage{fancyhdr}
\usepackage{titlesec}
\usepackage{booktabs}
\usepackage{amsmath,amssymb,amsfonts}
\usepackage{microtype}
\usepackage{xcolor}
\usepackage{mdframed}
\usepackage{parskip}
\usepackage{enumitem}
\usepackage{hyperref}

\definecolor{rulered}{RGB}{180,30,30}
\definecolor{boxgray}{RGB}{245,245,245}
\definecolor{keywordgray}{RGB}{100,100,100}

\pagestyle{fancy}
\fancyhf{}
\fancyhead[L]{\textsc{\small Claude Mag}}
\fancyhead[C]{\small\textit{Machine Learning Research Digest}}
\fancyhead[R]{\small 2026-11-04}
\fancyfoot[C]{\small\thepage}
\renewcommand{\headrulewidth}{0.8pt}

\titleformat{\section}{\normalsize\bfseries\color{rulered}}{}{0pt}{}%
  [\vspace{-4pt}\color{rulered}\rule{\columnwidth}{0.4pt}]
\titlespacing{\section}{0pt}{8pt}{4pt}

\hypersetup{colorlinks=true,urlcolor=rulered,linkcolor=black}

\begin{document}
\setlength{\parindent}{0pt}
\setlength{\parskip}{4pt}

{\small\color{keywordgray}\textbf{Keywords:}
  zeroth-order optimization \textbullet{}
  memory-efficient fine-tuning \textbullet{}
  variational multiscale methods \textbullet{}
  LLM adaptation}\par\vspace{4pt}

{\LARGE\bfseries\raggedright
  ENGINE: Variational MultiScale Zeroth-Order LLM Fine-Tuning}\par\vspace{4pt}

{\large\itshape\raggedright
  A Closed-Form Subgrid Correction from Computational Mechanics Yields
  2.61$\times$ Faster ZO Convergence at Zero Memory Overhead}\par\vspace{6pt}

{\small\textbf{By} Zhuoli Ouyang, Bowen Li, Chen Hao, Yejin Choi}\par\vspace{2pt}

{\small\color{keywordgray}arXiv:2611.03293 \textbullet{}
  cs.LG \textbullet{} Submitted 2 Nov 2026}
\par\vspace{2pt}\noindent{\color{rulered}\rule{\columnwidth}{0.6pt}}\vspace{6pt}

ENGINE applies the Variational MultiScale (VMS) decomposition from
computational fluid mechanics to zeroth-order LLM fine-tuning. By
decomposing perturbations into a low-rank coarse component and computing
a closed-form subgrid correction that reintroduces unresolved curvature
at no additional cost, ENGINE achieves a 2.61$\times$ speedup in forward
passes over prior ZO methods with zero additional peak memory, and
consistently outperforms state-of-the-art ZO baselines on NLP benchmarks.

\begin{mdframed}[backgroundcolor=boxgray,linecolor=rulered,linewidth=1pt,
    innerleftmargin=6pt,innerrightmargin=6pt,
    innertopmargin=6pt,innerbottommargin=6pt]
\textbf{\small AT A GLANCE}\par\vspace{4pt}
\begin{description}[leftmargin=0pt,labelsep=4pt,itemsep=2pt,parsep=0pt,
    font=\small\bfseries]
  \item[Problem:] Zeroth-order optimization for LLM fine-tuning is slow
    to converge because low-rank perturbation subspaces miss curvature
    from the complement space.
  \item[Why it matters:] ZO fine-tuning is the only way to train LLMs
    with RL rewards or in memory-limited settings without backpropagation;
    improving its efficiency makes it practical.
  \item[Method:] VMS decomposes the parameter space into a low-rank
    coarse component (updated with 2 forward passes) and a fine-scale
    complement; a closed-form subgrid correction reintroduces complement
    curvature with no extra forward passes.
  \item[Result:] 2.61$\times$ speedup in forward passes over LOREN;
    zero additional peak memory overhead; consistent improvements on
    SuperGLUE, MMLU, and instruction-following.
  \item[Evidence:] GPT-J-6B on SuperGLUE, Llama-3-8B on MMLU, and
    instruction-following on Alpaca Eval.
\end{description}
\end{mdframed}

\section{The Big Picture}

Zeroth-order (ZO) optimization enables fine-tuning LLMs without
backpropagation, reducing memory from $O(\text{parameters})$ to
$O(1)$ for the optimizer state. This matters in three settings: (1)
fine-tuning on GPUs with insufficient VRAM for gradients, (2) fine-tuning
through black-box API access where gradients are unavailable, and (3)
fine-tuning with non-differentiable reward signals such as human
feedback or unit-test pass rates.

The main limitation is convergence speed: estimating gradients by
measuring reward differences between perturbed parameters is noisy.
Restricting perturbations to a low-rank subspace (as in MeZO and LOREN)
reduces noise but introduces a systematic bias: curvature from the
complement of the subspace is ignored, slowing convergence.

ENGINE addresses this with an insight from computational mechanics:
the unresolved fine-scale curvature can be reintroduced via a closed-form
correction derived from the coarse (low-rank) gradient estimate, without
any additional forward passes.

\section{Why Existing Approaches Fall Short}

\textbf{Full-space ZO (MeZO)} estimates gradients with random directions
in the full parameter space. The variance is $O(d)$ where $d$ is the
parameter count, making it impractically noisy for billion-scale models.

\textbf{Low-rank ZO (LOREN)} restricts perturbations to a rank-$r$
subspace. This reduces variance to $O(r)$ but discards curvature from
the $d - r$ dimensional complement, introducing a systematic bias that
slows convergence. With $r \ll d$, the complement curvature is large.

\textbf{Increasing the rank $r$} reduces bias but increases memory and
noise proportionally. There is no free way to access the complement
curvature within the two-forward-pass budget.

\section{Core Mechanism: Variational MultiScale Decomposition}

The VMS framework decomposes the parameter space into a coarse
subspace $\mathcal{U}$ (spanned by a low-rank basis
$U \in \mathbb{R}^{d \times r}$) and its complement
$\mathcal{U}^\perp$. Any perturbation $z$ decomposes as:
\[
  z = z_c + z_f, \quad z_c = UU^\top z,\quad z_f = (I - UU^\top) z
\]

ENGINE estimates the coarse gradient using the standard two-point
ZO estimator restricted to $z_c$:
\[
  \hat{g}_c = \frac{\ell(\theta + \mu z_c) - \ell(\theta - \mu z_c)}{2\mu} \cdot z_c
\]
using 2 forward passes. The \textbf{subgrid correction}
$\delta_f$ accounts for the fine-scale contribution:
\[
  \delta_f = (I - UU^\top) H_\text{diag}\, \hat{g}_c
\]
where $H_\text{diag}$ is a diagonal approximation of the Hessian
maintained as an exponential moving average of squared gradient
magnitudes. This term has zero additional forward-pass cost and $O(d)$
memory, equal to the existing optimizer state.

The final update is:
\[
  \theta \leftarrow \theta - \eta\bigl(\hat{g}_c + \alpha\, \delta_f\bigr)
\]
where $\alpha$ is a correction weighting hyperparameter.

\section{Learning or Inference Procedure}

\textbf{Initialization:} Compute $U$ as the top-$r$ right singular
vectors of a batch of initial ZO gradient estimates. Maintain
$H_\text{diag}$ as an exponential moving average of $\hat{g}_c^2$.

\textbf{Per-step update:}
\begin{enumerate}[itemsep=2pt,parsep=0pt]
  \item Sample $z_c \sim \mathcal{N}(0, UU^\top)$ (a low-rank direction).
  \item Evaluate $\ell(\theta \pm \mu z_c)$ with 2 forward passes.
  \item Compute $\hat{g}_c$ and subgrid correction $\delta_f$.
  \item Update $\theta$ using $\hat{g}_c + \alpha\,\delta_f$.
  \item Update $H_\text{diag}$ with the new squared gradient components.
\end{enumerate}

\textbf{Periodic basis refresh:} Every $K = 200$ steps, recompute $U$
from recent gradient history via a lightweight SVD of a $d \times K$
matrix of gradient directions.

\section{What the Guarantee Says}

Proposition 1 (informal): The expected bias of the ENGINE gradient
estimate is strictly lower than that of LOREN with the same rank $r$,
provided $\alpha > 0$ and $H_\text{diag}$ approximates the diagonal
Hessian. No exact convergence rate is derived, but the proposition
bounds the bias reduction as a function of the quality of the diagonal
Hessian approximation.

The practical speed-up guarantee is empirical, not asymptotic. No
formal convergence rate is provided in the paper.

\section{Experimental Findings}

\begin{itemize}[itemsep=2pt,parsep=0pt]
  \item \textbf{SuperGLUE (GPT-J-6B):} ENGINE matches or exceeds
    full-fine-tuning accuracy on 7 of 8 tasks; achieves the same
    accuracy as LOREN with \textbf{2.61$\times$ fewer forward passes}.
  \item \textbf{MMLU (Llama-3-8B):} ENGINE achieves 72.3\%
    5-shot accuracy, vs.\ 69.1\% for LOREN and 70.8\% for MeZO at the
    same forward-pass budget.
  \item \textbf{Alpaca Eval:} ENGINE win-rate vs.\ LOREN is 61.4\%,
    confirming quality improvements in open-ended instruction following.
  \item \textbf{Memory:} Peak memory overhead of ENGINE vs.\ MeZO is
    0\%---the diagonal Hessian state fits within the existing optimizer
    state memory.
\end{itemize}

\section{Ablations and Interpretation}

\textbf{Without subgrid correction ($\alpha = 0$):} Performance equals
LOREN, confirming that the speedup is entirely due to the correction.

\textbf{Rank $r$ sensitivity:} ENGINE with $r = 32$ outperforms LOREN
with $r = 128$ on SuperGLUE, demonstrating that the correction
compensates for a fourfold reduction in rank.

\textbf{Hessian approximation quality:} Using a richer block-diagonal
Hessian estimate improves convergence by 7\% but requires 1.3$\times$
more memory. The diagonal approximation provides most of the benefit.

\textbf{Basis refresh interval $K$:} Performance is robust to $K$ in
the range 100--500, suggesting that the low-rank basis does not need
to be updated frequently.

\section*{Reference}

Zhuoli Ouyang, Bowen Li, Chen Hao, Yejin Choi.
\textbf{ENGINE: Endogenous Variational MultiScale Optimization for
Zeroth-Order LLM Fine-Tuning.}
arXiv:2611.03293. 2026.
\url{https://arxiv.org/abs/2611.03293}

\end{document}
